\documentclass[11pt]{article}

\usepackage[margin=25mm]{geometry}

\usepackage{amsmath,amssymb,fontspec,unicode-math}

\IfFontExistsTF{DejaVu Serif}{\setmainfont{DejaVu Serif}}{\setmainfont{Latin Modern Roman}}

\IfFontExistsTF{Latin Modern Math}{\setmathfont{Latin Modern Math}}{\setmathfont{DejaVu Math TeX Gyre}}

\usepackage{xurl,hyperref}

\setlength{\parindent}{0pt}\setlength{\parskip}{6pt}

\begin{document}

\title{VidEmo: Affective-Tree Reasoning for Emotion-Centric Video Foundation Models}

\author{ELVIA preliminary AI analysis}\date{21 October 2026}\maketitle

Preliminary assessment: Minor.\par

Independent scientific review remains pending. Author adoption has not been requested.

The current evaluation release can lose the actual labeled question and aggregate failed judge requests with inconsistent denominators. Original published score effects remain unestablished.

\section*{Review scope}

Full supplied technical report read alongside the corresponding original-paper claim passages and applicable author-release source. This pass independently checks the stated reasoning and bounds; historical executions were not rerun.

\section*{Problem 1: Scorer can replace the labeled question with a system instruction}

Source: Appendix D.3–D.4; eval/eval\_results.py, get\_response, lines 190–201

Assessment: supported; impact: minor

\subsection*{Source claim (paraphrase)}

The evaluator scores against messages[0] instead of selecting the user question by role. Released predictions include a system message before the question.

\subsection*{Analysis}

Consider two items with identical system instruction, reference letter A and model response happy. In one user question A means happy; in the other A means sad. Correct judgments differ, but selecting only the system message makes both judge prompts identical and omits both option lists. No deterministic or stochastic judge can recover the missing distinction from identical inputs. A directly inspected released VidEmo emotion record uses this system-first schema.

\subsection*{Finding and scope}

The supplied scorer and raw records have an information-losing interface mismatch. A private user-first cleaning stage might have been used historically; the published score impact is not established.

\subsection*{Proposed correction}

Select and validate the actual user message, preserve task instructions/options and required dialogue context, and remove only known video placeholders.

The described correction remains a proposal; author adoption or a corresponding erratum is not established.

\section*{Problem 2: Scoring failures change denominators inconsistently}

Source: Appendix D.3; eval/eval\_results.py, eval\_single and result\_evaluate, lines 215–269

Assessment: supported; impact: minor

\subsection*{Source claim (paraphrase)}

Exhausted judge requests can be absent from the stored results. Numeric scores average stored responses, while yes/no accuracy divides by all submitted records.

\subsection*{Analysis}

For two submitted items with only one successful score of \(100\), the numeric branch reports \(100\), while one successful yes gives accuracy \(50\). The missing numeric value could be anywhere in \([0,100]\), so the full mean is only identified within \([50,100]\). In general, if N are expected, m are scored and observed sum is A, the full mean lies in \([A/N,(A+100(N-m))/N]\).

\subsection*{Finding and scope}

The current metrics use different failure semantics and may summarize different populations. No actual original judge-failure rate or direction of benchmark bias is established; complete successful scoring avoids this problem.

\subsection*{Proposed correction}

Keep one indexed status for every expected record, validate accepted scores, and report complete means or an explicit failure policy. Conditional means and worst-case missing-score bounds should be labeled separately.

The described correction remains a proposal; author adoption or a corresponding erratum is not established.

\section*{Limits}

Preliminary AI source/argument assessment; independent human scientific adjudication remains pending.\par

No author-provided or retained verification program was executed. Published empirical results were not independently reproduced.\par

\section*{Source identity}

Authors’ arXiv:2511.02712v1, current official version checked 2026-10-21

\url{https://arxiv.org/pdf/2511.02712v1}

Reviewed PDF SHA-256: \nolinkurl{104baa156db4eb28f392b888273fd72ab51ed5f77db5fdbc11490e7e1fdef3d1}

The SHA-256 identifies the exact current-author-version PDF inspected. The original acquisition timestamp for this retained PDF is unknown; file modification times were not used. Current version identity was independently checked against official arXiv metadata.

\end{document}
